\chapter{Background}\label{chapter:background}

The chapter introduces the concepts needed to understand the system described in
the report and reviews research on multi-agent debate, sycophancy, and agent
influence. It then explains the gap addressed by the study.


\section{Preliminaries}\label{sec:preliminaries}
There are several techniques for developing "trust" in an Artificial Intelligence (AI) system. Some of these approaches depend upon a consensus of answers generated by multiple AI systems. Other techniques depend upon finding relevant scientific information and assessing the accuracy of the information found.
Chain-of-Thought prompts assist in illustrating the reasoning process that a model uses to arrive at its conclusion, so that we may better understand how an answer was produced and why an error was made during an instance of poor performance~\cite{wei2022chain,zhang2025from}. A Multi-Agent Debate Framework provides a methodology where each agent evaluates the arguments presented by the other agents. Each agent presents their proposed solutions and includes support for those proposed solutions. All agents evaluate the proposed solutions and supporting rationales provided by the other agents. Following evaluation and revision of their respective positions and solutions, a final proposed solution is selected~\cite{du2024improving}. One advantage of allowing each agent to independently evaluate a problem is that if each agent discovers different faults in the system, then collectively the entire team may identify potential problems with the system. Additionally, a team member can utilize another team member's rationalization to criticize aspects of their own rationalization.
However, regardless of how well the individual components of the solution are developed, the combination of results (i.e., the methodology of combining results) is equally as important as the results themselves. Most commonly, results are combined utilizing a majority vote. In a majority vote, every member has an equal say, and the solution supported by the highest number of members is therefore selected. However, majority voting does not necessarily eliminate errors; rather, it can strengthen repeated errors~\cite{wang-etal-2024-rethinking-bounds}. Sycophantic Consensus occurs when an agent agrees with the majority position or changes a correct position in favor of a more strongly supported but incorrect position~\cite{pitre-etal-2025-consensagent}.
There are two types of confidence associated with voting: Self-Confidence in one's position and Token Distribution Confidence. Neither type of confidence verifies the truth claims contained within a position~\cite{yoffe-etal-2025-debunc}. DebUnc identifies this weakness by comparing uncertainty metrics produced by the deployed system with a Diagnostic Ground Truth Metric that already contains the known answer~\cite{yoffe-etal-2025-debunc}. Developing trust in an AI system therefore requires obtaining and evaluating scientifically relevant data and evidence.


\section{Literature Review}\label{sec:literature-review}
The review categorizes the reviewed literature by five categories: debate construction, external verification, sycophancy, uncertainty, and debate evaluation. This study primarily compares methods for determining when to generate a response (i.e., methods for deciding on debate), methods for changing how an agent interacts with another agent (i.e., methods for changing agent interaction), and the method for changing what influences an agent's decision based upon information retrieved from outside the system (i.e., changing agent influence). Recent surveys are provided in~\cite{li2024survey} as they outline the design space of multi-agent systems utilizing LLMs. In order to develop generated responses that are more reliable, research has shifted focus toward structured reasoning, multi-agent collaboration, and external validation of responses generated by LLMs~\cite{zhang2025from}. Multi-Agent Debate is an example of this new area of research because Multi-Agent Debate allows for other agents to evaluate and potentially challenge each agent's reasoning before that agent produces their final output~\cite{du2024improving}. The next section outlines some of the most relevant current works in this emerging area and illustrates how the proposed architecture addresses an existing research gap.

\subsection{Similar Applications}\label{subsec:similar-applications}
It should be noted that this system is very different from a standard web/mobility application. In particular, the system we are building is an experimental layer designed to enable deliberative processes and act as an integration layer for scientific question answering systems, research assistants, evidence based decision support tools and/or other applications requiring reliable reasoning. As opposed to the traditional approach where one model solution is generated, the systems listed above generate multiple candidate solutions. Additionally, the systems listed above allow for interaction among candidates. Finally, each system aggregates results from the candidates to generate a final answer for users. Du et al. showed that structured debate among several agents increases factuality and reasoning compared to single-agent approaches~\cite{du2024improving}. Wang et al. showed that layered aggregation can result in higher overall model capability when combining outputs from several models without supporting explicit debate~\cite{wang2025mixture}. Multi-agents systems differ primarily regarding how output is combined. Some systems utilize structured debate as a means to enhance reasoning between agents~\cite{du2024improving}, while others use hierarchical pipelines to aggregate/average the output from each pipeline (the individual agents)~\cite{wang2025mixture}. Zhou et al. proposed a self-organizing system for tasks to determine whether a specialized reasoner produces the response versus being in line with a pre-defined workflow~\cite{zhou2025reso}. Additionally, Kushwaha \& Madhavan further developed an augmented retrieval framework that identifies contradictions found within retrieved documents and ranks credibility of the source prior to producing the response~\cite{kushwaha2026agentbased}. Both are examples of how distributed reasoning (across multiple specialized agents) results in increased robustness/accuracy in factually correct responses. What matters is on what scale you construct dialogue. Ueda et al. analyzed how to construct multi-agent LLM dialogues to facilitate research ideation~\cite{ueda2025exploring}; however, Wang et al. studied if/how large-scale multi-agent LLMs could provide assistance with scientific code generation~\cite{wang2025scalability}. Their aggregation method does not change during dialogue. Similar to a single LLM, the proposed framework has the same reliability goals for answers; yet allows for various levels of influence for each agent throughout the discussion process. After each round of dialogue, evidence verification results are generated; and after each round, the trust scores for all agents are updated based on claim-level scientific evidence supporting their claims. Therefore, an unsupported majority consensus would hold less value than a well-supported minority view.

\subsection{Related Research}
\label{subsec:related-research}

The work of Wei et al. (2022) shows that ``Chain-of-Thought Prompting''
improves the ability of a large language model to reason~\cite{wei2022chain}.
The work of Du et al. (2024) extended the work of Wei et al. (2022) by
incorporating multi-agent debate in addition to chain-of-thought prompting,
in which multiple agents give answers to a question, look at how other agents
have reasoned when answering the question, and then adjust their reasoning
~\cite{du2024improving}.In contrast to both previous approaches for enhancing the performance of
large language models, CRITIC~\cite{gou2024critic} relies on outside tools
to validate the answers provided and enhance them. It therefore permits
validation of assertions contained in each response; however, there is no
mechanism to compare additional models working on the same question
simultaneously. Instead, Mixture-of-Agents develops numerous independent
agent-proposer models that produce answers to identical questions and layers
the responses into one output, indicating what several large language models
can do together; however, the layering of the individual agent responses is
fixed and contains no evidence-based trust values~\cite{wang2025mixture}.
iMAD addresses cost rather than correctness~\cite{fan2026imad}. Fan, Yoon,
and Ji identified 41 properties of a structurally defined self-critic and
constructed a classifier that determines whether or not to initiate a debate
based on those properties. They were able to attain token reductions of up to
92\% relative to GroupDebate and increases in accuracy of up to 13.5\% over
a single-agent baseline. Though agent influence was constant following the
initiation of the debate, agent influence remained constant throughout the
debate.Pitre, Ramakrishnan, and Wang investigated sycophancy in debates
~\cite{pitre-etal-2025-consensagent}. The results of their investigation
indicate that the correct answer existed in more than 20\% of cases where
the agents responded incorrectly. CONSENSAGENT identifies stalling and
answer copying in a debate. When either condition is identified, CONSENSAGENT
alters the task prompt before initiating the next stage of the debate. The
results of the study demonstrate that measured sycophancy decreased as a
result; however, the total trust score is still dependent upon self-reported
confidence and consistency.DebUnc permits agents to express uncertainty among themselves~\cite{yoffe-etal-2025-debunc}. However, agents express uncertainty via prompts or attention scaling. The deployed measures provided very limited benefit. A diagnostic ground-truth signal is available for comparison purposes. The diagnostic ground-truth signal has access to information concerning the correct answer and has greater potential than the deployed measures. The difference between the deployed measures and the diagnostic ground-truth signal indicates that the major barrier is the quality of the trust signal.Kushwaha and Madhavan created an 8-agent MAS-RAG system intended to identify contradictions between retrieved documents and rank evidence based upon credibility~\cite{kushwaha2026agentbased}. Although source credibility was evaluated by humans prior to deliberation, it was never changed during deliberation.
There have also been numerous studies evaluating the social aspects of
debate. The authors of Cau et al. identified how selective persuasion could
cause an increase in opinion convergence~\cite{cau2025selective}. Yeow et al.
studied decision architectures and found that feedback between agents can
spread early errors~\cite{yeow2025reliable}. Wang et al. demonstrated how a
good example given in context within one large language model (LLM) may
perform better than multi-agent dialogue in certain situations. As noted in
~\cite{wang-etal-2024-rethinking-bounds}, judge error and peer conformity
were identified as potential drawbacks.
The authors of Choi, Zhu, and Li separated communication used in debate from
the act of voting and found that majority voting explained a substantial
portion of the normal performance gains~\cite{choi2025debate}. FREE-MAD
extended previous research by eliminating the need for final agreement and
utilizing anti-conformity prompts and trajectory scores~\cite{cui2026freemad}.
ReSo utilizes dynamic task graphs with reward-based processes for choosing
agents~\cite{zhou2025reso}.
Even though past research has explored optimizing coordination and the choice
of agents, none of the previously conducted studies have assessed scientific
claims prior to altering the influence of an agent.
Sycophancy also occurs at the model level. The results of training models to
be perceived as ``warm'' and ``agreeable,'' as demonstrated by Ibrahim et al.,
produced lower levels of factual accuracy and higher levels of sycophancy
~\cite{ibrahim2026training}. Chen et al. utilized self-augmented preference
alignment from a training perspective to address the issue
~\cite{chen2025self}.
Process-level diagnostics for engagement, responsiveness, influence
asymmetry, balance, stability, and agent utility developed by Pitre et al.
would be helpful when ground truth is unavailable
~\cite{pitre2026diagnostic}. Final agreement is insufficient to measure a
debate, as stated by Pitre et al.~\cite{pitre2026diagnostic}.


\section{Gap Analysis}
\label{sec:gap-analysis}
There are still three open issues with regard to Multi-Agent Debate Systems (MADS) from our review of the current literature. The first issue is whether the extra computational burden associated with a MADS can be justified. Fan et al. have used iMAD for this purpose -- to predict whether a debate is needed prior to the start of a debate~\cite{fan2026imad}. The second issue is how to detect unsupported agreement caused by conformity or sycophancy. CONSENSAGENT addresses this problem through prompt optimization during the debate~\cite{pitre-etal-2025-consensagent}. The third issue is how to determine which agents should receive greater influence when they disagree with one another. DebUnc has a weighting mechanism to deal with uncertainty as it relates to the debate process; however, the trust mechanisms of this weighting rely on each individual agent's internal assessment of uncertainty and therefore do not use external validation or independently verifiable evidence~\cite{yoffe-etal-2025-debunc}.As previously mentioned, even when multiple agents agree with respect to a particular piece of information, that information may be incorrect. For example, Cau et al. showed that selective persuasion can affect the level of consensus reached by different agents~\cite{cau2025selective}, while Wang et al. showed that peer conformity can limit the effectiveness of multi-agent debate~\cite{wang-etal-2024-rethinking-bounds}. While existing retrieval-based systems either verify the correctness of retrieved information or attempt to improve the aggregation process, these systems do not dynamically update an agent's influence based upon verified scientific evidence during a debate~\cite{gou2024critic}. Similarly, although dynamic coordination mechanisms help solve the issue of selecting which agents participate in the debate, they do not provide a mechanism for increasing or decreasing an agent's influence based on evidence of its behavior during the debate~\cite{zhou2025reso}. The existing research therefore does not currently contain a method by which an agent's influence can be continuously updated using independently verified evidence of its behavior in real time.The proposed framework addresses this remaining gap by decomposing each agent's responses into individual factual claims and identifying scientific passages that relate to those claims. Each claim is then classified as either supported, contradicted, or unverifiable. The classification of each claim serves as input to an update function for a bounded trust score. As a result, well-supported minority agents can continue to have influence during a debate, while an unsupported majority consensus is given less weight. To evaluate the proposed framework, we will use a number of metrics, including Consensus-Collapse Rate (CCR), Minority-Preservation Rate (MPR), Evidence-Calibration Rate (ECR), and Answer Accuracy (AA), alongside conventional answer accuracy measures.



\section{Summary}\label{sec:background-summary}

The chapter introduced LLM reasoning, chain-of-thought prompting, multi-agent
debate, and sycophantic consensus. The reviewed studies cover debate efficiency,
prompt correction, uncertainty weighting, external verification, process evaluation,
and agent coordination. They do not continuously update agent influence from
claim-level scientific evidence. The next chapter uses this gap to define the
system requirements, detailed methodology, and project plan.
